%! Simplifying Rational Expressions — Unit 4, Lesson 4.1 — Student Packet Cover Sheet
\documentclass[10pt]{article}
\usepackage{algebra2-article}
\usepackage{algebra2-boxes}
\usepackage{ltablex}
\keepXColumns

\begin{document}

% Full-bleed forest banner
\begin{tikzpicture}[remember picture, overlay]
  \fill[forest] (current page.north west)
    rectangle ([yshift=-0.9in]current page.north east);
\end{tikzpicture}
\vspace{-0.8in}

\begin{center}
  {\color{white}\textbf{\LARGE Algebra 2}}\\[4pt]
  {\color{forestmid}\large\textbf{Unit 4 \quad Rational Functions}}\\[6pt]
  {\color{white}\textbf{\large Lesson 4.1 \quad Simplifying Rational Expressions \& Domain Restrictions}}
\end{center}

\vspace{0.22in}
\namedateperiod
\vspace{0.18in}

\begin{learningtargetbox}
\small
\begin{enumerate}[leftmargin=1.5em, itemsep=5pt]
  \item \ldots{} \textbf{factor} a rational expression completely and write it in \textbf{simplest form} by dividing out common \textbf{factors} --- never terms.
  \item \ldots{} list \textbf{every domain restriction} from the \textbf{original} denominator, including the ones whose factor cancels, and say which produce a \textbf{hole} and which a \textbf{vertical asymptote}.
  \item \ldots{} \textbf{justify} that two rational expressions are \textbf{equivalent}, and name exactly which inputs they disagree on.
\end{enumerate}
\end{learningtargetbox}

\vspace{0.14in}

\begin{tocbox}
\small
\renewcommand{\arraystretch}{1.5}
\begin{tabularx}{\linewidth}{@{} c l X r @{}}
\rowcolor{forestbg}
\textbf{\#} & \textbf{Component} & \textbf{Description} & \textbf{Score} \\
\midrule
1 & Warm-Up        & Spiral review: factoring completely, factors vs.\ terms, and finding excluded values & \blank{1.2cm} \\
2 & Guided Notes   & The four steps, opposite binomials, and what cancelling does \emph{not} change & \blank{1.2cm} \\
3 & Group Activity & \emph{Factor, Restrict, Divide Out} --- a feature table, an error analysis, and equivalence, by tier & \blank{1.2cm} \\
4 & Exit Ticket    & Quick check: simplify with restrictions, plus an SOL-style opposite-binomial item & \blank{1.2cm} \\
5 & Homework       & Monomial and binomial practice, a graph with a hole, and a design-your-own extension & \blank{1.2cm} \\
\midrule
  & \multicolumn{2}{r}{\textbf{Total}} & \blank{1.2cm} \\
\end{tabularx}
\end{tocbox}

\vspace{0.10in}

\begin{remindbox}
\small A \textbf{rational expression} is a fraction of polynomials, and simplifying one is four steps:
\textbf{factor} the top and bottom completely, \textbf{list the restrictions} from the
\emph{original} denominator, \textbf{divide out} the shared factors, and \textbf{write the simplest
form with those restrictions attached}. Cancelling is legal because $\frac{c}{c}=1$ --- which is why
you may only divide out a common \textbf{factor} (something multiplied), never a \textbf{term}
(something added): $\frac{x+3}{3}\neq x$, and one test value proves it. Two expressions are
\textbf{equivalent} when they agree at every input \emph{both} accept --- so simplifying never
restores a thrown-out input. That is the whole point: a restriction whose factor \emph{cancels}
becomes a \textbf{hole}, and one whose factor \emph{stays} becomes a \textbf{vertical asymptote} ---
the two features you read off graphs in Lesson 4.0, now produced by algebra. Watch for
\textbf{opposite binomials} such as $5-x$ and $x-5$: they do not cancel, they divide to $-1$.
\end{remindbox}

\end{document}
